\documentclass[acmsmall,screen]{acmart}

\begin{document}

\acmConference[POPL '26]{Principles of Programming Languages}{January 11--17}{Rennes, France}

\begin{abstract}
  Extensible cases, dual to row polymorphism, and
\end{abstract}

\title{Extensible Cases Plus Pattern Match as Inverse Realizes Linear Logic's De Morgan Laws and Enforce Exhaustiveness for Free}
\author{V. E. McHale}
\orcid{0000-0001-6093-2967}
\affiliation{
  \institution{Northern Trust}
  \city{Chicago}
  \country{United States}
}
\email{vamchale@gmail.com}
\maketitle

\section{Background}

Extensible cases were implemented by Blume, Acar, and Chae \cite{blume2006}, solving the expression problem.

% linear logic: two inverses, "fruitful interaction"

\section{Realizing Linear Logic}

Munch-Maccagnoni \cite{munchmaccagnoni2009} points out that linear logic's $\&$ ("with") can be the type of pattern matching.

\bibliographystyle{ACM-Reference-Format}
\bibliography{pm.bib}

\end{document}
